\section{曲线命名}

给曲线起个名字对读者很有帮助。以由下式定义的函数 $f$ 为例：

\[
   x>0\ \text{且}\ f(x)=\dfrac{x^2+x-1}{x^3}
\]

用 \tkzcname{tkzFnText} 宏很容易写出标题。坐标以坐标轴单位给出。要沿曲线放置文字，最简单的方法是定义曲线上的一个点并使用其名称。

\begin{tkzltxexample}[num]
  \tkzFnDefPoint(f,3){P}
  \tkzFnText[above right](P){${\mathcal{C}}_f$}
\end{tkzltxexample}

第一行定义图像上 $x=3$ 处的点 $P$。第二行在 $P$ 处放置文字。\TIKZ\ 选项可调整效果。

\subsection{添加标签}

\begin{center}
\begin{tkzexample}[vbox]
\begin{tikzpicture}
  \tkzFnInit[xmin=0,xmax=10,
             ymin=0,ymax=1.2,ystep=0.2]
  \tkzFnGrid
  \tkzFnAxes
  \tkzFnDraw[name=f,thick,color=red,domain=0.55:10]{(x*x+x-1)/x^3}
  \tkzFnText(3,-0.3){\textbf{$\mathbf{f}$ 的图像}}
  \tkzFnDefPoint(f,3){P}
  \tkzFnText[above right,text=red](P){${\mathcal{C}}_f$}
\end{tikzpicture}
\end{tkzexample}
\end{center}

\endinput
